\documentclass{article}
\usepackage[margin=1.5cm]{geometry}
\usepackage{amsmath}
\usepackage{enumitem}

\setlist{nosep,leftmargin=*}

\begin{document}
\twocolumn

\title{Problem Set 4: Functions and Modules}
\author{Prof. Jordan C. Hanson}
\date{}

\maketitle

\section{Exercise 1: Patient Data Module}

\begin{enumerate}

\item Create a Python module called

\begin{verbatim}
patient_data.py
\end{verbatim}

The module will organize a block of fictional patient records.
Each patient should be represented by a tuple containing

\begin{verbatim}
(name, ssn, age, dob,
 height, weight, gender)
\end{verbatim}

Use the following data:

\begin{verbatim}
patients = [
 ("Alan Reed", "000-00-0001", 74,
  "04/12/1952", 70, 218, "M"),

 ("Maria Lopez", "000-00-0002", 76,
  "09/03/1950", 64, 155, "F"),

 ("Robert Chen", "000-00-0003", 68,
  "01/19/1958", 72, 225, "M"),

 ("James Hill", "000-00-0004", 82,
  "11/28/1943", 69, 210, "M"),

 ("Susan Patel", "000-00-0005", 72,
  "06/14/1954", 65, 205, "F"),

 ("David King", "000-00-0006", 71,
  "02/07/1955", 71, 195, "M"),

 ("Henry Cole", "000-00-0007", 79,
  "08/21/1947", 68, 242, "M"),

 ("Linda Grant", "000-00-0008", 66,
  "12/02/1959", 63, 180, "F")
]
\end{verbatim}

The SSNs above are fictional identifiers created only for this
programming exercise.

\end{enumerate}


\section{Exercise 2: Module Functions}

\begin{enumerate}

\item Inside \texttt{patient\_data.py}, define the following
functions.

\begin{itemize}

\item A function that returns the complete patient list:

\begin{verbatim}
def get_patients():
    """Return all patient records."""
    return patients
\end{verbatim}

\item A function called

\begin{verbatim}
def screening_rule(patient):
\end{verbatim}

that returns \texttt{True} only when all three conditions are met:

\[
\text{gender = male},
\qquad
\text{age} > 70,
\qquad
\text{weight} > 200.
\]

Use the tuple indices and Boolean operator \texttt{and}.
For example,

\begin{verbatim}
age = patient[2]
weight = patient[5]
gender = patient[6]
\end{verbatim}

The function must \textbf{return} a Boolean value rather than
printing the result.

\item Define a third function

\begin{verbatim}
def display_patient(patient):
\end{verbatim}

that prints the patient's information in a readable format.

For example,

\begin{verbatim}
Name: Alan Reed
SSN: 000-00-0001
Age: 74
DoB: 04/12/1952
Height: 70 in
Weight: 218 lb
Gender: M
\end{verbatim}

Include a short docstring inside each function.

\end{itemize}
\end{enumerate}


\section{Exercise 3: Screening Program}

\begin{enumerate}

\item Create a second Python file called

\begin{verbatim}
heart_screen.py
\end{verbatim}

Import your module using

\begin{verbatim}
import patient_data
\end{verbatim}

Retrieve the patient records:

\begin{verbatim}
patients = patient_data.get_patients()
\end{verbatim}

Then use a \texttt{for} loop to examine each patient:

\begin{verbatim}
for patient in patients:

    if patient_data.screening_rule(
            patient):

        patient_data.display_patient(
            patient)

        print("")
\end{verbatim}

Your program should display \textbf{only} patients who satisfy
all three screening conditions.

\begin{itemize}

\item Do not repeat the screening logic in
\texttt{heart\_screen.py}.  The decision must be made by the
function inside the module.

\item Do not manually type the names of qualifying patients.

\item Verify that changing one patient's age, weight, or gender in
\texttt{patient\_data.py} can change whether that patient appears
in the final output.

\item Add the heading

\begin{verbatim}
Patients Matching Screening Rule
\end{verbatim}

before displaying the selected records.

\end{itemize}

\end{enumerate}


\section{Check Your Program}

\begin{enumerate}

\item Your submission should contain exactly two Python files:

\begin{verbatim}
patient_data.py
heart_screen.py
\end{verbatim}

The module should store the patient data and define reusable
functions.  The main program should import the module and call
those functions.

For this programming exercise, the three-condition rule is only a
software-filtering example and should not be interpreted as an
actual medical assessment of heart-disease risk.

\end{enumerate}

\end{document}